\subsection{European Natural Gas Markets}\label{sec:gas-markets}
Natural gas is a physical commodity that exists within a gas grid. The Title Transfer Facility (TTF) is a virtual trading hub in the Netherlands. A virtual trading hub allows participants to transfer ownership without reference to a specific physical location. TTF is operated by Gasunie Transport Services (GTS) within the Dutch gas transmission network.  Gas futures contracts referencing TTF are traded on ICE Endex, the European energy exchange operated by Intercontinental Exchange. Futures contracts are described in Sec.~\ref{sec:futures}.

A market is considered liquid if many participants can buy or sell continuously, without affecting the asset price. TTF is the most liquid natural gas market in Europe, with traded volumes 4.5 times the combined volume of the other eight main European hubs \parencite{heatherEuropeanTradedGas2026}. Therefore, the TTF gas price serves as the reference price for gas contracts. 

At virtual hubs such as TTF, gas price formation is determined by short-run supply and demand \parencite{obadiDRIVINGFUNDAMENTALSNATURAL2020}. In the short run the infrastructure of the gas grid is fixed, and the supply of gas cannot expand in response to a change in demand. Therefore, the adjustment takes place through the price, rather than the quantity.

Gas demand originates in three sectors. The first source is residential and commercial heating, which makes consumption seasonal.
Deviations of weather from the seasonal norm also affect prices \parencite{brownWhatDrivesNatural2008}. Industrial production is the second source of demand. In the chemical sector, for example, gas is used as an energy source and as a raw material. The third source is power generation, where gas competes with coal. As a result, the gas market is linked to the coal market and, through emission allowances, to the carbon market.

On the supply side, natural gas prices depend on pipeline flows and deliveries of liquefied natural gas (LNG). LNG is natural gas cooled to liquid form for transport by ship. 

As has been observed repeatedly in the past, the European natural gas market is highly sensitive to supply disruptions. The European gas price crisis between 2021 and 2022 is the most severe episode. During the COVID-19 pandemic in 2020, energy demand collapsed. This was followed by economies reopening, and gas demand recovered faster than production could follow. From mid-2021, the post-pandemic increase in demand was also combined with reduced Russian pipeline flows, strong Asian competition for LNG cargoes and storage levels below the seasonal norm \parencite{heather2022unfortunate}. As a result, TTF gas prices rose from roughly €18/MWh in the first quarter of 2021 to a peak of €339/MWh on 26 August 2022 and then declined from early 2023. 

The most recent disruption followed the military campaign against Iran by Israel and the United States on 28 February 2026. Iran responded with attacks on shipping and energy infrastructure in the Gulf. Between 27 February and 2 March, the traffic through the Strait of Hormuz fell from 95 vessels to 4 vessels per day \parencite{persian-gulf-26}. These events resulted in a spike of TTF gas prices from €27/MWh on 5 January 2026 to €61/MWh on 19 March 2026.

The European gas price crisis eventually raised the question of if the TTF price benchmark had become less representative of the underlying gas market. In November 2022, the European Commission proposed a Market Correction Mechanism (MCM), which would cap gas prices during crisis periods \parencite{MCM-euro-commision}. The mechanism applied until January 2025 but was never activated, as gas prices declined after 2023.

\subsection{Futures Contracts}\label{sec:futures}
Futures contracts are standardized agreements to buy or sell an asset at
a fixed price on a specified future date. 
The quantity and the delivery date are set by the exchange. The price is agreed between the parties and remains fixed for the contract. As a result, all contracts for a given month are identical in their terms and can be freely transferred. TTF gas futures are quoted in EUR/MWh and are physically delivered, with delivery spread evenly across the hours of the delivery month. 

Price formation of natural gas in Europe has generally taken place on the futures rather than spot market \parencite{obadiDRIVINGFUNDAMENTALSNATURAL2020}. A futures contract can be traded repeatedly before maturity. As a result, the market is open for participants who hedge or speculate without an interest in the physical commodity.

At any point in time, contracts for multiple consecutive months are listed and can be traded in parallel. The contract with the nearest delivery on a given day is the front-month contract (M1), and the next is the second-month contract (M2). Trading concentrates in the front-month contract, hence it is commonly used as the reference for the current gas price.

\subsection{Options and Implied Volatility}\label{sec:options}
\todo{
    1. what is an option, call/put, strike, maturity, payoff, valuation (P vs Q, brief)
    2. Black-76 and its assumptions, implied volatility, volatility smile
    3. Closing link: smile shape reflects skewness and kurtosis, connecting to Sec. 2.4;
       option prices depend on the distribution of the underlying at maturity, which is
       what makes them a test of a generative model
}


\subsection{Stylized Facts of Commodity Returns}
A diverse collection of assets, markets and time periods share common statistical characteristics of returns. These characteristics are referred to as stylized facts. Eleven stylized facts are summarized in \textcite{contEmpiricalPropertiesAsset2001}. 

\todo{As described in Sec.~\ref{sec:options}, the value of a European option depends on the
distribution of the underlying price at maturity, and in particular on the probability of
large moves away from the current level. A model that generates return paths for option
pricing must therefore reproduce the shape of the return distribution and not only its
mean. write this whn 2.3 is written 
alternaitve: The payoff of a European option is determined by the price of the underlying
at maturity. Its value therefore depends on the distribution of that price,
and in particular on the probability of large moves away from the current
level. }
Only three stylized facts are measured for TTF in this thesis: heavy tails, volatility clustering and asymmetry. 

\textbf{Heavy tails.} The tails of a distribution are the regions far from its mean. Large and infrequent returns are located in the tails. The distribution of financial returns has more mass in the tails than a normal distribution \parencite{contEmpiricalPropertiesAsset2001}. This is referred to as heavy tails. The weight of the tails is measured by the excess kurtosis

\begin{equation}
    \mathrm{Kurt}_{\mathrm{ex}}(X) = \mathbb{E}\!\left[\left(\frac{X-\mu}{\sigma}\right)^{4}\right] - 3,
    \label{eq:kurtosis}
\end{equation}
where $\mu$ and $\sigma$ are the mean and the standard deviation of the random variable $X$. The excess kurtosis of a normal distribution is 0. The distribution is considered heavy-tailed if its excess kurtosis is positive. \todo{Heavy tails determine the value of out-of-the-money
options, which pay nothing unless a large move occurs.}

\textbf{Volatility clustering.} Large returns are more likely to be followed by large returns, and small returns by small returns  \parencite{contEmpiricalPropertiesAsset2001}. This is measured with the autocorrelation. Autocorrelation is the correlation of a time series with itself, shifted by a fixed number
of observations called the lag. The autocorrelation of returns at lag $\tau$ is 
\begin{equation}
    \rho(\tau) = \mathrm{corr}\!\left( r_{t+\tau},\ r_{t} \right),
    \label{eq:autocorrelation}
\end{equation}
where $r_t$ is the return on trading day $t$. The returns themselves show almost no autocorrelation, while absolute or squared returns are positively correlated \parencite{contEmpiricalPropertiesAsset2001}. Volatility clustering is measured on squared returns
\begin{equation}
    C_2(\tau) = \mathrm{corr}\!\left( r_{t+\tau}^{2},\ r_{t}^{2} \right),
    \label{eq:vol-clustering}
\end{equation}
$C_2$ remains positive through squaring and decays slowly over several days.
\todo{Volatility clustering determines how option value scales with time to maturity. Under
independent returns, the aggregated return over several days approaches a normal
distribution and the tails thin out. Clustering slows this, so that tail weight is preserved
at longer horizons.}

\textbf{Gain/loss Asymmetry.}

\todo{6. Scope exclusion. Mean reversion and seasonality are commodity stylized facts you don't measure — say why. Daily log returns of a rolled front-month series difference away most of the price-level mean reversion, and seasonality lives in the term structure rather than in daily returns. Two or three sentences, cite Schwartz (1997).  Forward pointer. Measured for TTF in 4.1, for generated returns in Chapter 4.}